% Class-compatibility preamble retained for a reader excerpt of
%   Lisa Sauermann, Dmitrii Zakharov, "A Sharp Ramsey Theorem for Ordered
%   Hypergraph Matchings", arXiv:2309.04813v4, CC BY 4.0.
%
% The source pins the publisher class aic (the AIMS journal class), which is not part of
% the delivery baseline, so this file supplies the amsart baseline and then copies the
% author's own declaration block mechanically and unchanged
% (pinned-original-source.tex lines 25 and 28-59: the ams packages, the author's eleven
% newtheorem declarations with [theorem] as the shared counter, the definition theorem
% style, the author's comment that the file should add as few macro or package
% definitions as possible, and the author's own macros \su, \sm, \Z, \HH, \conv and \PP).
% The author's tikz package (pinned source line 26) serves the three figures of the paper
% that this excerpt does not carry and is therefore not reproduced; the author's babel,
% fontenc and inputenc lines (pinned source lines 47-49) are layout/font choices of the
% publisher build and are not part of the mathematical excerpt, the delivery engine being
% XeLaTeX with UTF-8 input.
% No mathematical macro, statement, symbol or assumption is introduced, changed or
% removed by this file.
\documentclass[10pt,a4paper]{amsart}
\usepackage{amsthm,amsfonts,amssymb,amsmath}
\usepackage{graphicx}

\newtheorem{theorem}{Theorem}[section]
\newtheorem{proposition}[theorem]{Proposition}
\newtheorem{lemma}[theorem]{Lemma}
\newtheorem{conjecture}[theorem]{Conjecture}
\newtheorem{fact}[theorem]{Fact}
\newtheorem{claim}[theorem]{Claim}
\newtheorem{corollary}[theorem]{Corollary}
\newtheorem{definition}[theorem]{Definition}
\newtheorem{problem}[theorem]{Problem}
\newtheorem{question}[theorem]{Question}
\newtheorem{observation}[theorem]{Observation}

\theoremstyle{definition}
\newtheorem{remark}[theorem]{Remark}

% The author notes at this point that as few macro or package
% definitions as possible should be added; the six macros below are the
% author's own and are copied unchanged.

\newcommand{\su}{\subseteq}
\newcommand{\sm}{\setminus}
\newcommand{\Z}{\mathbb{Z}}
\newcommand{\HH}{\mathcal{H}}
\newcommand{\conv}{\operatorname{conv}}
\newcommand{\PP}{\mathcal{P}}
\usepackage{cleveref}
\crefname{originalequation}{original equation}{original equations}
\Crefname{originalequation}{Original equation}{Original equations}
\makeatletter
\newlabel{Lp}{{4.1}{source}{}{sourcepointer-1}{}}
\newlabel{Lp@cref}{{[originalequation][0][]4.1}{[1][0][]source}{}{}{}}
\makeatother
\crefname{originallemma}{original lemma}{original lemmas}
\Crefname{originallemma}{Original lemma}{Original lemmas}
\makeatletter
\newlabel{obs-construction-partite}{{4.2}{source}{}{sourcepointer-2}{}}
\newlabel{obs-construction-partite@cref}{{[originallemma][0][]4.2}{[1][0][]source}{}{}{}}
\makeatother
\title{A sharp lower-bound construction for ordered hypergraph matchings along a fixed chain of ordered partitions}
\author{Lisa Sauermann, Dmitrii Zakharov}
\date{arxiv:2309.04813v4; CC BY 4.0}
\begin{document}
\maketitle
\noindent\textit{Editorial note.} All delivered mathematical text is the author's own, copied from the pinned author source sources/combinatorics-fl02-a-sharp-ramsey-theorem-for-ordered/pinned-original-source.tex; no statement or proof step was rewritten, reordered or removed.\par\medskip
\noindent\textit{Editorial note.} Preamble provenance: the source pins the publisher class aic, which is not part of the delivery baseline, so this card uses sources/combinatorics-fl02-a-sharp-ramsey-theorem-for-ordered/retained-preamble.tex, which supplies the amsart baseline and then copies the author's ams packages, the author's eleven newtheorem declarations sharing the theorem counter, the definition theorem style and the author's own macros \textbackslash{}su, \textbackslash{}sm, \textbackslash{}Z, \textbackslash{}HH, \textbackslash{}conv and \textbackslash{}PP (pinned source lines 25 and 28-59) verbatim. No macro was invented and no mathematical symbol was changed.\par\medskip
\noindent\textit{Editorial note.} Bibliography provenance: the citations of the counted proof are resolved with sources/combinatorics-fl02-a-sharp-ramsey-theorem-for-ordered/bibliography.tex, an excerpt of the author's own thebibliography environment (pinned source lines 598-703) copied character for character with all eighteen entries kept; no entry was rewritten, reordered, merged, split or invented and no citation key was changed.\par\medskip
\noindent\textit{Editorial note.} External references are navigation pointers only; every deliberately kept step of the author's argument stays inside the counted spans. Their printed numbers are the author's own numbering under the author's newtheorem declarations, computed mechanically by the lane's scanner.\par\medskip
\noindent\textit{Editorial note.} The labels the counted proof invokes outside its own span, Lp and obs-construction-partite, are navigation-pointer stubs; the construction lemma has its own card in this lane (P1-ramsey-partite-construction) and is declared here as a named standard prerequisite.\par\medskip
\noindent\textit{Allowed standard prerequisites (editorial).}\begin{itemize}
\item The construction of an r-partite ordered hypergraph matching with prescribed clique-size bounds. The author's lemma at source lines 556-560, whose proof at source lines 564-571 is delivered by this lane's companion card for that statement. Supplies the building block that the theorem iterates along the chain of ordered partitions.
\item The quantity L\_P of a collection and its role in the clique-size bounds. Defined in the author's preliminaries (source lines 225-281) and used through the display labelled Lp at source line 549. Names the quantities that the construction keeps below the prescribed bounds.
\end{itemize}\medskip
\ifdefined\hypertarget\hypertarget{sourcepointer-1}{}\fi
\noindent\textit{Source reference.} Original equation 4.1: sources/combinatorics-fl02-a-sharp-ramsey-theorem-for-ordered/pinned-original-source.tex, source lines 549--549. Navigation pointer only: it labels the author's display defining the quantity L\_P; the display belongs to the surrounding section and is not reproduced in this card.\par\medskip
\ifdefined\hypertarget\hypertarget{sourcepointer-2}{}\fi
\noindent\textit{Source reference.} Original lemma 4.2: sources/combinatorics-fl02-a-sharp-ramsey-theorem-for-ordered/pinned-original-source.tex, source lines 556--556. Navigation pointer only: the counted proof iterates this construction lemma; it has its own card in this lane and its text is declared here as a named standard prerequisite.\par\medskip
\begin{theorem} \label{thm-paramteters-lower-bound-stronger}
Let $r$ be a positive integer, and let $\lambda^{(1)} \succ \lambda^{(2)} \succ \dots \succ \lambda^{(r)}$ be a sequence of ordered partitions of $r$. For every $r$-pattern $P\in \PP(\lambda^{(1)})\cup\dots\cup \PP(\lambda^{(r)})$, consider a positive integer $m_P$. Then there exists an  ordered $r$-uniform hypergraph matching $\HH$ of size
\[|\HH|=\prod_{P\in \PP(\lambda^{(1)})\cup\dots\cup \PP(\lambda^{(r)})} m_{P},\]
such that any two edges $e,f\in \HH$ form an $r$-pattern in $\PP(\lambda^{(1)})\cup\dots\cup \PP(\lambda^{(r)})$, and for each $r$-pattern $P\in \PP(\lambda^{(1)})\cup\dots\cup \PP(\lambda^{(r)})$ every $P$-clique contained in $\HH$ has size at most $m_P$.
\end{theorem}

\begin{proof}[Proof of Theorem~\ref{thm-paramteters-lower-bound-stronger}]
We prove the theorem by induction on $r$. Note that the base case $r=1$ is trivial, so let us assume $r\ge 2$ and that we already proved the theorem for $r-1$.

In the sequence $\lambda^{(1)} \succ \lambda^{(2)} \succ \dots \succ \lambda^{(r)}$, each ordered partition $\lambda^{(s)}$ for $s=1,\dots,r$ must have precisely $s$ parts. In particular, the partition $\lambda^{(r-1)}$ has the form $\lambda^{(r-1)}=(1,\dots,1,2,1,\dots,1)$. Let $j\in \{1,\dots,r-1\}$ be the unique index such that $\lambda^{(r-1)}_j=2$. If we imagine $\lambda^{(1)} \succ \lambda^{(2)} \succ \dots \succ \lambda^{(r)}$ as partitions of the set $\{1,\dots,r\}$ into intervals, then the partition  $\lambda^{(r-1)}$ consists of the set $\{j,j+1\}$ and singleton sets for all the remaining elements. Hence $j$ and $j+1$ are in the same set for each of the partitions $\lambda^{(1)} \succ \lambda^{(2)} \succ \dots \succ \lambda^{(r-1)}$. If we ignore $j$ and relabel the remaining elements accordingly (i.e.\ $j+1$ is relabeled $j$, $j+2$ is relabeled $j+1$, and so on), $\lambda^{(1)} \succ \lambda^{(2)} \succ \dots \succ \lambda^{(r-1)}$ turns into a sequence of ordered partitions $\mu^{(1)} \succ \mu^{(2)} \succ \dots \succ \mu^{(r-1)}$ of $r-1$. For $s=1,\dots,r-1$ we obtain $\lambda^{(j)}$ from $\mu^{(j)}$ by doubling the element $j$ (keeping both copies in the same set of the partition), and relabeling the ground set into $\{1,\dots,r\}$ accordingly.

Now, for $s=1,\dots,r-1$, every $r$-pattern $P\in \PP(\lambda^{(s)})$ turns into an $(r-1)$-pattern $P'\in \PP(\mu^{(s)})$ when omitting the $j$-th $\rm A$ and the $j$-th $\rm B$ in $P$ (note that these occur in the same block, i.e.\ right adjacent to the $(j+1)$-th $\rm A$ and the $(j+1)$-th $\rm B$, respectively). Conversely, $P$ can be recovered from $P'\in \PP(\mu^{(s)})$ by doubling the $j$-th $\rm A$
in $P$ into $\rm{AA}$ and doubling the $j$-th $\rm B$
in $P$ into $\rm{BB}$. Thus, these operations give natural bijections between $\PP(\lambda^{(s)})$ and $\PP(\mu^{(s)})$.

For every $(r-1)$-pattern $P'\in \PP(\mu^{(1)})\cup\dots\cup \PP(\mu^{(r-1)})$, let us define $m_{P'}=m_P$, where $P\in \PP(\lambda^{(1)})\cup\dots\cup \PP(\lambda^{(r-1)})$ is the $r$-pattern obtained from $P'$ by doubling the $j$-th $\rm A$ and the $j$-th $\rm B$. Then, applying the induction hypothesis to the sequence $\mu^{(1)} \succ \mu^{(2)} \succ \dots \succ \mu^{(r-1)}$ of ordered partitions of $r-1$, we obtain an ordered $(r-1)$-uniform hypergraph matching $\HH'$ of size
\[|\HH'|=\prod_{P'\in \PP(\mu^{(1)})\cup\dots\cup \PP(\mu^{(r-1)})} m_{P'}=\prod_{P\in \PP(\lambda^{(1)})\cup\dots\cup \PP(\lambda^{(r-1)})} m_{P},\]
such that $L_{P'}(\HH')\le m_{P'}$ for all $(r-1)$-patterns $P'\in \PP(\mu^{(1)})\cup\dots\cup \PP(\mu^{(r-1)})$ and $L_{P'}(\HH')=1$ for all other $(r-1)$-patterns $P'\not\in \PP(\mu^{(1)})\cup\dots\cup \PP(\mu^{(r-1)})$.

From $\HH'$, we can now obtain an ordered $r$-uniform hypergraph matching $\HH^*$ with $|\HH^*|=|\HH'|$ by doubling the $j$-th element of every edge $e'\in \HH'$. More formally, for every edge $e'\in \HH'$, we add a new element to $e'$ right before the $j$-the element of $e'$ (in the given ordering of the underlying vertex set of $\HH'$). For the resulting edge $e\in \HH^*$, we can recover the original edge $e'\in \HH'$ by deleting the $j$-the element of $e$. If $e,f\in \HH^*$ form some $r$-pattern $P$, then the $(r-1)$-pattern formed by the corresponding edges $e',f'\in \HH'$ is the $(r-1)$-pattern $P'$ obtained from $P$ by omitting the $j$-th $\rm A$ and the $j$-th $\rm B$. Thus, by the properties of $\HH'$, we have $L_P(\HH)\le m_P$ for all $r$-patterns $P\in \PP(\lambda^{(1)})\cup\dots\cup \PP(\lambda^{(r-1)})$ and $L_P(\HH)=1$ for all other $r$-patterns $P\not\in \PP(\lambda^{(1)})\cup\dots\cup \PP(\lambda^{(r-1)})$.

Finally, recall that $\lambda^{(r)}=(1,\dots,1)$ is the unique ordered partition of $r$ into $r$ parts. Let $\HH_r$ be an $r$-partite ordered $r$-uniform hypergraph matching as in Lemma~\ref{obs-construction-partite}. Now, the blow-up $\HH=\HH^*[\HH_r]$ has size
\[|\HH|=|\HH^*|\cdot |\HH_r|=|\HH'|\cdot |\HH_r|=\prod_{P\in \PP(\lambda^{(1)})\cup\dots\cup \PP(\lambda^{(r-1)})} m_{P}\cdot \prod_{P\in \PP(\lambda^{(r)})} m_{P}=\prod_{P\in \PP(\lambda^{(1)})\cup\dots\cup \PP(\lambda^{(r)})} m_{P},\]
and by (\ref{Lp}) for every $r$-pattern $P$ we have $L_P(\HH)\le L_P(\HH^*)\cdot L_P(\HH_r)$. Hence $L_P(\HH)\le m_P\cdot 1=m_P$ for all $P\in \PP(\lambda^{(1)})\cup\dots\cup \PP(\lambda^{(r-1)})$, and $L_P(\HH)\le 1\cdot m_P=m_P$ for all $P\in \PP(\lambda^{(r)})$, and also $L_P(\HH)\le 1\cdot 1=1$ for all $P\not\in \PP(\lambda^{(1)})\cup\dots\cup \PP(\lambda^{(r)})$. Thus, we have constructed the desired ordered $r$-uniform hypergraph matching $\HH$.
\end{proof}

% Bibliography excerpt retained for reader excerpts of
%   Lisa Sauermann, Dmitrii Zakharov, "A Sharp Ramsey Theorem for Ordered Hypergraph Matchings",
%   arXiv:2309.04813v4, CC BY 4.0.
%
% Every entry below is copied character for character from the author's own thebibliography
% environment of the pinned author source (lines 598-703).  No entry was rewritten, reordered,
% merged, split or invented and no citation key was changed; all 18 entries are kept so that
% any citation of a delivered excerpt resolves without guessing an entry.
\begin{thebibliography}{99}

%\bibitem{bergelson-johnson-moreira}
%Vitaly Bergelson, John H. Johnson Jr., and Joel Moreira.
%\newblock New polynomial and multidimensional extensions of classical partition
%  results.
%\newblock 2015, arXiv:1501.02408.

%\bibitem{cilleruelo}
%Javier Cilleruelo.
%\newblock Combinatorial problems in finite fields and {S}idon sets.
%\newblock {\em Combinatorica}, 32(5):497--511, 2012.


\bibitem{AJKS} 
Michael Anastos, Zhihan Jin, Matthew Kwan, and Benny Sudakov.
\newblock Extremal, enumerative and probabilistic results on ordered hypergraph matchings.
\newblock {\em Forum of Mathematics Sigma} 13: Paper No.~e55, 39~pp, 2025.

\bibitem{Campos} 
Marcelo Campos, Simon Griffiths, Robert Morris, and Julian Sahasrabudhe. 
\newblock An exponential improvement for diagonal Ramsey.
\newblock 2023, arXiv:2303.09521.

\bibitem{Caro} 
Yair Caro. 
\newblock New Results on the Independence Number.
\newblock Technical Report, Tel-Aviv University, 1979.

\bibitem{ExpES} 
Recep Altar Çiçeksiz, Recep Altar, Zhihan Jin, Eero Räty, and István Tomon. 
\newblock Exponential Erd\H{o}s-Szekeres theorem for matrices.
\newblock {\em Proceedings of the London Mathematical Society} 129, no. 3: Paper No. e12632, 14~pp, 2024.

\bibitem{Dil} 
Robert P. Dilworth 
\newblock A Decomposition Theorem for Partially Ordered Sets.
\newblock {\em Annals of Mathematics}  51: 161--166, 1950.

\bibitem{DGR22}
Andrzej Dudek, Jarosław Grytczuk, and Andrzej Ruci\'{n}ski. 
\newblock Ordered unavoidable sub-structures in matchings and random matchings.
\newblock {\em Electronic Journal of Combinatorics} 31, no. 2: Paper No. 2.15, 27 pp, 2024

\bibitem{DGR} 
Andrzej Dudek, Jarosław Grytczuk, and Andrzej Ruci\'{n}ski. 
\newblock Erd\H{o}s-Szekeres type Theorems for ordered uniform matchings.
\newblock {\em Journal of Combinatorial Theory, Series B} 170: 225--259, 2025.

\bibitem{ES} 
Paul Erd{\H o}s and George Szekeres. 
\newblock A combinatorial problem in geometry.
\newblock {\em Compositio Mathematica} 2: 463--470, 1935.

\bibitem{GL} Timothy W. Gowers and Jason Long.
\newblock The length of an s-increasing sequence of r-tuples. \newblock {\em Combinatorics, Probability and Computing} 30, no. 5: 686--721, 2021.

\bibitem{graham-rothschild-spencer} 
Ronald L. Graham, Bruce L. Rothschild, and Joel H. Spencer. 
\newblock {\em Ramsey theory.} Vol. 20. John Wiley \& Sons, 1991.


\bibitem{Lee} 
Choongbum Lee. 
\newblock Ramsey numbers of degenerate graphs.
\newblock {\em Annals of Mathematics} 185, no. 3: 791--829, 2017.

\bibitem{LS} 
Nathan Linial and Michael Simkin. 
\newblock Monotone subsequences in high-dimensional permutations. 
\newblock {\em Combinatorics, Probability and Computing} 27, no. 1: 69--83, 2018.

\bibitem{MV} 
Sam Mattheus and Jacques Verstraete. 
\newblock The asymptotics of $r(4,t)$.
\newblock {\em Annals of Mathematics} 199, no.~2: 919–941,  2024.

%\bibitem{Mirsky} 
%Leon Mirsky. 
%\newblock A dual of Dilworth’s decomposition theorem.
%\newblock {\em Amer. Math. Monthly} 78, no. 8: 876-877, 1971.

\bibitem{MS} 
Guy Moshkovitz and Asaf Shapira. 
\newblock Ramsey theory, integer partitions and a new proof of the Erdős–Szekeres theorem. 
\newblock {\em Advances in Mathematics} 262: 1107--1129, 2014.

\bibitem{promel} 
Hans Jürgen Prömel. 
\newblock {\em Ramsey theory for discrete structures.} Vol. 3. New York: Springer, 2013.

\bibitem{RR} 
Christian Reiher, and Vojt\v{e}ch R\"odl. 
\newblock The girth Ramsey theorem.
\newblock 2023, arXiv:2308.15589.

\bibitem{ST} 
Tibor Szab\'{o} and G\'{a}bor Tardos. 
\newblock A multidimensional generalization of the Erd\H{o}s–Szekeres lemma on monotone subsequences.
\newblock {\em Combinatorics, Probability and Computing} 10, no. 6: 557--565, 2001.

\bibitem{Wei} 
Victor K. Wei. 
\newblock A Lower Bound on the Stability Number of a Simple Graph.
\newblock {\em Technical memorandum}, TM 81-11217-9, Bell Laboratories, 1981.


\end{thebibliography}
\end{document}
